\documentclass[10pt]{article}
\usepackage[margin=1in]{geometry}
\usepackage[T1]{fontenc}

\title{Generalized Tic Tac Toe Agent Report}
\author{}
\date{}

\begin{document}
\maketitle
\thispagestyle{empty}

\section*{Project Overview}
This project implements an adversarial-search agent for generalized Tic Tac Toe on an $n \times n$ board with target length $m$. The same search core is used in two settings: a local AI-vs-AI simulator and an online runner that connects to the NotExponential game API. The objective is to produce moves quickly enough for live games while still searching deeply enough to detect wins, blocks, and tactical traps.

The main design choice in this project is to avoid rebuilding the full game state after every move. Instead, the board stores incremental information about windows of length $m$, frontier cells near existing pieces, and a Zobrist hash for state lookup. This lets the agent spend most of its time on search rather than repeated board recomputation.

\section*{Evaluation Function}
The evaluation function combines \emph{line potential} and \emph{center control}. For every contiguous row, column, or diagonal window of length $m$, the implementation keeps track of how many X and O symbols appear in that window. If a window contains both players, it is treated as blocked and contributes zero. Otherwise, the score depends only on the number of pieces owned by one side in that open window.

The scoring weights grow exponentially, so stronger threats receive much larger values than weak patterns. In particular, windows with $m-1$ pieces are boosted heavily because they represent immediate winning chances or urgent defensive responsibilities. This makes the heuristic much more tactical than a simple piece-count function. It distinguishes clearly between a harmless dispersed position and a nearly completed line.

Center control is added as a secondary term. Cells closer to the geometric center are given higher weight, because central placements usually participate in more potential winning windows and create more follow-up options. This term is tapered as the board fills up, so it has strong influence in the opening but gradually gives way to direct tactical considerations in the middle and late game. In effect, the evaluation prefers central development early, then switches attention toward concrete threats once the board becomes crowded.

This heuristic is a good fit for generalized Tic Tac Toe because the game is dominated by line creation, line blocking, and move locality. It is also inexpensive to update incrementally: after a move, only the windows touching that cell need to be adjusted.

\section*{Minimax / Adversarial Search}
The agent uses negamax with alpha-beta pruning. Negamax is a compact form of minimax that takes advantage of the zero-sum structure of the game: the value for one player is the negation of the value for the opponent. Alpha-beta pruning reduces the number of explored nodes by cutting off branches that cannot improve the current result.

Search is wrapped inside iterative deepening. The agent first searches depth 1, then depth 2, and so on until it reaches the configured maximum depth or the move deadline. This has two benefits. First, the program always has a legal move available even under a short time budget. Second, the best move found at a shallower depth can be reused as a preferred move for deeper searches, which improves move ordering and therefore improves pruning.

Terminal states are handled explicitly. Wins and losses are assigned values near a large constant, and the score is adjusted by ply so that faster wins are preferred and delayed losses are preferred over immediate ones. If the board is full without a winner, the state is evaluated as a draw.

\section*{Performance Improvements}
Several implementation choices were added to make search practical on larger boards:

\begin{itemize}
\item \textbf{Incremental board state.} The board caches window counts, total line balance, center balance, and the winner state. A move and undo operation updates these values directly rather than scanning the whole board.
\item \textbf{Frontier-based candidate generation.} Instead of considering every empty cell, the move generator focuses on cells within a small radius of existing pieces. This is important for large boards, where most distant empty cells are strategically irrelevant.
\item \textbf{Aggressive move ordering.} Candidate moves are classified by tactical priority: immediate wins, immediate blocks, fork creation, fork prevention, defense, attack, and center gain. Better ordering leads to earlier alpha-beta cutoffs.
\item \textbf{Candidate limits by ply.} The root, second ply, and deeper plies use different candidate limits. The search remains broad where move quality matters most and becomes more selective deeper in the tree.
\item \textbf{Transposition table.} A Zobrist hash identifies board states, and a transposition table stores previously computed search results with exact, lower-bound, or upper-bound flags. This reduces repeated work when the same state is reached by different move orders.
\item \textbf{Time management.} The search uses a deadline with a safety margin and periodic timeout checks, ensuring the agent returns a move reliably during interactive play.
\end{itemize}

Two small tactical shortcuts are also useful at the root: if the agent already has an immediate winning move, it returns that move directly; if there is exactly one opponent winning move to block, it blocks immediately. These checks are cheap and prevent wasted search effort on forced situations.

\section*{Validation and Final Remarks}
The codebase includes unit tests for both state handling and agent behavior. The tests verify that making and undoing a move restores cached state correctly, that rebuilding the board from API-style coordinates matches the incremental state, and that the agent chooses immediate wins, blocks forced losses, and still returns legal moves under very short time budgets. There is also a consistency test showing that the transposition-table version and the uncached version agree on the selected move for the same position.

Overall, the project emphasizes a practical balance between search quality and response time. The heuristic is intentionally tactical, the search algorithm is standard but effective, and the main strength of the implementation comes from engineering choices that shrink the branching factor and avoid redundant computation. That combination is appropriate for competitive generalized Tic Tac Toe, where the difference between a usable agent and an unusable one is often not the theoretical algorithm alone, but how efficiently the surrounding board logic supports it.

\end{document}
